% ------------------------------------------------------------------------------
% Formatvorlage für RFCCC1
% ------------------------------------------------------------------------------
%   erstellt von Christian Schiffler, 25.11.2010
%   basierend auf Arbeiten von Stefan Macke, 24.04.2009
%   http://blog.stefan-macke.de

% Dokumentenkopf ---------------------------------------------------------------
%   Diese Vorlage basiert auf "scrbook" aus dem koma-script.
% ------------------------------------------------------------------------------
\documentclass[
    fontsize=11pt, % Schriftgröße
    DIV10,
    a4paper, % Papierformat
    %oneside, % einseitiges Dokument
    openany, % remove empty pages after a chapter if chapter starts on the "left" hand
    titlepage, % es wird eine Titelseite verwendet
    parskip=half, % Abstand zwischen Absätzen (halbe Zeile)
    headings=normal, % Größe der Überschriften verkleinern
    listof=totoc, % Verzeichnisse im Inhaltsverzeichnis aufführen
    bibliography=totoc, % Literaturverzeichnis im Inhaltsverzeichnis aufführen
    index=totoc, % Index im Inhaltsverzeichnis aufführen
    captions=tableheading, % Beschriftung von Tabellen unterhalb ausgeben
    final % Status des Dokuments (final/draft)
]{scrbook}

% Umlaute ----------------------------------------------------------------------
%   Umlaute/Sonderzeichen wie äüöß direkt im Quelltext verwenden (CodePage).
%   Erlaubt automatische Trennung von Worten mit Umlauten.
% ------------------------------------------------------------------------------
\usepackage[utf8x]{inputenc}
\usepackage[T1]{fontenc}
\usepackage{textcomp} % Euro-Zeichen etc.

% Meta-Informationen -----------------------------------------------------------
%   Informationen über das Dokument, wie z.B. Titel, Autor, etc.
%   werden in der Datei Meta.tex definiert und können danach global
%   verwendet werden.
% ------------------------------------------------------------------------------
% Meta-Informationen -----------------------------------------------------------
%   Definition von globalen Parametern, die im gesamten Dokument verwendet
%   werden können (z.B auf dem Deckblatt etc.).
%
%   ACHTUNG: Wenn die Texte Umlaute oder ein Esszet enthalten, muss der folgende
%            Befehl bereits an dieser Stelle aktiviert werden:
%            \usepackage[latin1]{inputenc}
% ------------------------------------------------------------------------------
\newcommand{\titel}{Request for Contao Coder Comments 1 - RFCCC1}
\newcommand{\untertitel}{ Contao Extension Repository 3}
\newcommand{\autor}{Tristan Lins, Christian Schiffler}
\newcommand{\jahr}{2011}
\newcommand{\logo}{images/base/Title}
\newcommand{\deckblattlogo}{images/base/Title}


% benötigte Packages -----------------------------------------------------------
%   LaTeX-Packages, die benötigt werden, sind in die Datei Packages.tex
%   "ausgelagert", um diese Vorlage möglichst übersichtlich zu halten.
% ------------------------------------------------------------------------------
\chapter{Packages}

\section[sec:versioning]{Versioning}
The versioning is based on the classic version numbering scheme with extensions.

A version number is combined like the following.

\begin{lstlisting}[caption=Version number in EBNF]
DOT            = <the dot (".") character>
UPALPHA        = <any US-ASCII upper case letter "A".."Z">
LOALPHA        = <any US-ASCII lower case letter "a".."z">
DIGIT          = <any US-ASCII digit "0".."9">
ALPHA          = UPALPHA | LOALPHA
ALNUM          = ALPHA | DIGIT
MAJOR          = DIGIT
MINOR          = DIGIT
REVISION       = DIGIT
EXTENSION      = <"-"> *ALPHA *DIGIT
BUILD          = *DIGIT
HEADBUILD      = <"head"> [DOT *ALNUM]
REGULARBUILD   = <"alpha" DIGIT | "beta" DIGIT | "rc" DIGIT | "stable"> [DOT BUILD]
RELEASE        = HEADBUILD | REGULARBUILD

VERSION = MAJOR DOT MINOR [DOT REVISION] [EXTENSION] DOT RELEASE
\end{lstlisting}

The only mandatory fields in the version are major, minor and the release. 
The release number is combined of either the word "head" in conjunction with an alpha numeric revision suffix which refers to unstable releases or one of the words "alpha", "beta", "rc" or "stable" suffixed with one digit. The latter one can also be accompanied with a build number.
The revision number (also known as "maintenance" or "patch level") is entirely optional, as it is not used in most cases.

The optional extension component is mostly useful for personal repositories, which serve altered versions of certain extensions. This is known from derivative works like the Ubuntu project, which is releasing altered debian packages with a suffix "ubuntu1" or similar.

Please note that developers are encouraged to provide adjustments to upstream maintainers in favor of providing own versions of packages though to keep the forking of extensions to a minimum.

\subsection[sec:versioning real word examples]{A real world versioning example}

Let's imagine an agency has altered an extension for ACME Inc. and is providing this extension via their own repository.
To eliminate problems with the vanilla extension in the main repository, the agency appends the extension "acme" to the version and results in the version
\DisplFile{1.2.3-acme1}.

Using this scheme one can prevent unwanted upgrading (see \ref{sec:versioncompatibility}).

In contrast to the other components of the version, the build number is the only one that can not be influenced by the developer. The build number is incremented by one for every package being uploaded into the repository but not when the description or the properties are changing. This is in contrast to the old ER, which incremented the build number upon every change to the extension.
\todo{How shall this work? version bumping can only be done upon "building" of a package but the upload is much later in the work flow.}


\section[sec:versioncompatibility]{Version compatibility}
Thanks to the \techTerm{extension} component of the version, it is possible to tag modified variants of an extension. These variants are in most cases incompatible to newer (vanilla) releases of the same extension.

To address these potentional problems, the following version compatibility rules are in place.

Versions are compatible if, and only if, any of the following conditions is met:
\begin{itemize}
  \item ...the alphabetic representation of the \techTerm{extension} component is identical.
  \item ...the extension component is omitted on both versions.
\end{itemize}

\small
\begin{longtable}{|p{0.25\textwidth}|p{0.05\textwidth}|p{0.25\textwidth}|}
\caption{Version compatibility examples (compatible: "="; incompatible: "!")} \\
\hline
\label{tab:versioncompatexamples}
\textbf{Version 1} & \textbf{=/!} & \textbf{Version 2} \\
\hline
 1.2.3 & = & 4.5.6
\\ \hline
 1.2-extA.3 & = & 4.5-extA.6
\\ \hline
 1.2-extB.3 & ! & 4.5-extC.6
\\ \hline
 1.2-extD3.4 & = & 5.6-extD7.8
\\ \hline
 1.2-extE3.4 & ! & 5.6-extF7.8
\\ \hline
\end{longtable}
\normalsize

Please note that this only refers to the compatibility in respect to the version and that the compatibility can be further restricted by the extension meta file as described in \ref{sec:package information file}.

\section[sec:version comparison]{Version comparison}

When comparing compatible versions (according to the rules defined in \ref{sec:versioncompatibility}), one has to compare versions with extension component with a different approach than those without extension component.

\subsection[sec:version comparison without extension]{Version comparison without extension component}
A version without extension component will be matched according to the usual mechanism.

A given version A is bigger (aka. newer) than version B if, and only if, one of the following conditions is true:
\begin{itemize}
\item A.major>B.major
\item A.major=B.major \& A.minor>B.minor
\item A.major=B.major \& A.minor=B.minor \& A.revision>B.revision
\item A.major=B.major \& A.minor=B.minor \& A.revision=B.revision \& A.build>B.build
\end{itemize}

A given version A ist identical to version B if, and only if:
\begin{itemize}
\item A.major=B.major \& A.minor=B.minor \& A.revision=B.revision \& A.build=B.build
\end{itemize}

\subsection[sec:version comparison with extension]{Version comparison with extension component}

If an \techTerm{extension} component is present, the following additional rules are to be applied in addition to the rules mentioned in \ref{sec:version comparison without extension}.

A given version A is bigger (aka. newer) than version B if, and only if:
\begin{itemize}
\item A.extensionA = B.extensionA \& A.extensionN > B.extensionN
\end{itemize}

A given version A ist identical to version B if, and only if:
\begin{itemize}
\item A.extensionA = B.extensionA \& A.extensionN = B.extensionN
\end{itemize}

At this, one must pay attention to the fact that the alphabetical representation and the numeric representation has to be evaluated separately.

If the alphabetical representation (extensionA) is different for two versions, they have to be regarded as incompatible according to \ref{sec:versioncompatibility}. When no numeric representation (extensionN) is present in a version with extension, the version has to be evaluated as if a numeric part has been set with the value of "1".

\small
\begin{longtable}{|p{0.25\textwidth}|p{0.25\textwidth}|p{0.25\textwidth}|}
\caption{Version comparison example} \\
\hline
\label{tab:versioncompareexamples}
\textbf{field} & \textbf{Version 1} & \textbf{Version 2} \\
\hline
full version & 1.2-acme3.4 & 1.2-acme4.1
\\ \hline
major & 1 & 1
\\ \hline
minor & 2 & 2
\\ \hline
extension & acme3 & acme4
\\ \hline
extensionA (alpha) & acme & acme
\\ \hline
extensionN (numeric) & 3 & 4
\\ \hline
\end{longtable}
\normalsize

\section[sec:packet structure]{Package structure}
Extensions are shipped in a package file. Such a package contains meta information, optional package management routines for installation, un-installation and handling upgrades in addition to the application data.

The package file is a zip file with a changed extension to mark it as a package file. This approach is already known \ie from java (jar files).

For the file extension currently the following are considered (on the right hand are the file formats listed that are known to use this file extension):
\\
\begin{tabular}{lp{0.9\textwidth}}\\
.cext & none known so far\\
.cex  & INMOS Transputer Development System Occam User Program, ThumbsPlus File, The Currency Exchanger Rate File, CLAN Output File (Child Language Data Exchange System)\\
.ctp  & American Greetings CreataCard (Broderbund), BestAddress HTML Editor Combined Template File (Multimedia Australia Pty. Ltd.)\\
.ctx  & CTX Chemical File (Gasteiger Group), Chinese Character Input File, Compressed Text, GE Industrial Systems CIMPLICITY Text Version HMI Screen, Alphacam Compiled Text (Planit), Online Course Text (Microsoft Corporation), Pretty Good Privacy (PGP) Ciphertext File (PGP Corporation), Windows Terminal Server INI Backup File (Microsoft Corporation), Visual Basic User Control Binary File, CTRAN/W DEFINE Compressed Data File (Geo-Slope International), Nokia PC Suite Backup Contact File (Nokia)\\
\end{tabular}

A package contains the following files:
\begin{itemize}
\item PACKAGE - package information (required, see \ref{sec:package information file})
\item FILES - list of files contained in the package (required for certification, see \ref{sec:files information file})
\item LICENSE files - the license files for the extension (required according to \ref{sec:package information file})
\item ... any files the developer might need below the virtual root "CONTAO" which refers to the contao installation root.
\item ... any files the developer might need below the virtual root "DATA" for the user data root.
\end{itemize}

\subsection[sec:package information file]{PACKAGE information file}

The XML format is described in the schema file extension.xsd\supref{lst:extension.xsd}.

As usage example serves extension.xml\supref{lst:extension.xml}.

\subsubsection[sec:dependencies conflicts suggestions]{Package relationships - dependencies, conflicts and preferations}

A relationship of an extension to another extension can be defined in form of a \techTerm{dependency}, a \techTerm{conflict} or a \techTerm{preferation}.
Each of these relationships can be towards any version of the mentioned extension or even be narrowed down to certain version numbers.

\subsubsection[sec:versioning of relationships]{Versioning of relationships}

As dependencies and conflicts are often regarding certain versions or a certain range of versions.

In the definition of an relationship, we may use the following attributes:
\begin{itemize}
\item \techTerm{match} which means an exact match on the version (X=Y).
\item \techTerm{lowest} which is implicit inclusive and tells the lower end of the range (X>=Y).
\item \techTerm{highest} which is implicit exclusive and tells the upper end of the range (X<Y).
\end{itemize}
The attributes \techTerm{lowest} and \techTerm{highest} do not have to be used in combination, therefore it is possible to define an open range.
When used in combination, they define a closed range.

The attributes \techTerm{lowest} and \techTerm{highest} must not be used in conjunction with the attribute \techTerm{match}. If used in conjunction with \techTerm{match}, the exact match takes precedence over the other attributes and therefore is always exact.

In the package information file, this might look like the following when defining a conflict for a certain db driver (in this case mysql).
\lstset{language=XML} 
\begin{lstlisting}[caption=XML example: conflict example mysql]
...
<ctx:conflict name="db:mysql">
    <ctx:version match="4"/>
    <ctx:version lowest="5" highest="5.1"/>
    <ctx:version lowest="5.2"/>
</ctx:conflict>
...
\end{lstlisting}

In above example definition we have a conflict for:
\begin{itemize}
\item all version 4 drivers.
\item all versions 5 to 5.1 ( 5<=X<5.1)
\item all versions from 5.2 on upwards (5.2<=X)
\end{itemize}
With above condition we define a conflict for all versions of MySQL except for version 5.1 (we could have done it more simple though).

\subsubsection[sec:handling of versions with extension component]{Handling of versions with extension component}

In general the extension component of a version is ignored as it makes common sense that an altered version of an extension is compatible to it's unaltered counterpart from an external point of view. This behaviour can be altered by defining the extension version in the relationship explicitly. If specified, the match for the version must be performed including the extension part.

\subsubsection[sec:special relationships]{Special relationships}

Package relationships can not only refer other packages but some other name-spaces as well.
This is useful to restrict dependency on some php features, a certain database driver, some PEAR package or something else that has been defined in the specification.

Currently the following name spaces exist.\\
\begin{tabular}{ll}\\
contao         & The Contao core.\\
php:[feature]  & Relates to a certain php feature. See php features in \ref{sec:package relationship name spaces php features}\\
pear:[package] & Relates to a certain pear package. See pear packages in \ref{sec:package relationship name spaces pear packages}.\\
db:[type]      & Relates to a certain database, \ie db:mysql. See database drivers in \ref{sec:package relationship name spaces database drivers}.\\
\end{tabular}

\subsubsection[sec:alternative dependencies]{Alternative dependencies}
Often it is required to have a certain feature available but multiple ways exist to get this feature available when depending on some PHP feature but to also provide a wrapper class as workaround if it is not available.

Instead of defining features within the repository client and therefore having to update the client all the time when a new feature has been defined, we provide the developer with the possibility to define alternative paths by specifying groups of dependencies.

When defining groups of dependencies, the first positive match found is accepted and the group evaluation will be stopped. Note that the order in the XML has no meaning and therefore the attribute \techTerm{preferation} must be used to define which dependency is the most preferable.

\lstset{language=XML} 
\begin{lstlisting}[caption=XML example: alternative relationships]
...
<ctx:group preferation="php:soap nusoap">
<ctx:dependency name="php:soap">
    <ctx:description>I love soap, thats because i need this php extension.</ctx:description>
</ctx:dependency>
<ctx:dependency name="nusoap">
    <ctx:description>If you dont have soap, I pleased with an alternative implementation.</ctx:description>
</ctx:dependency>
</ctx:group>
...
\end{lstlisting}
This example describes, that it depends on the PHP SOAP extension but alternatively can also use the extension "nusoap" if the first is not available.

\subsubsection[sec:upgrade and downgrade compatibility]{Upgrade and downgrade compatibility}
In many cases upgrades and downgrades may cause problems. This is even more the case, when the data schema is changing.
That's the reason why the developer needs a way to define what potential threat an upgrade or downgrade of a certain package implies.

This potentional threat can be categorized in the following sections:

\begin{tabular}{llp{0.55\textwidth}}\\
\textbf{Category} & \textbf{threat factor} & \textbf{Description}\\
lossless  & none & Up- and downgrade can be performed without data loss.\\
dropUnused  & almost none & Up- and downgrade will remove data that will not be in use any more.\\
partial  & might be dangerous & Up- and downgrade might remove user data. Therefore manual intervention is suggested.\\
undefined & dragons live here! &This setting is implicit defined when the developer has not defined a setting at all. The result might be user data loss, in worst case scenario resulting in total data loss.\\
\end{tabular}

Whenever a developer has not defined a compatibility, the following automatic assumptions are made:\\
When upgrading, the compatibility is defaulting to "partial".
When downgrading, the compatibility is defaulting to "undefined".
The extension repository manager will then display a message where the up/downgrade compatibility is shown in addition to the description text of the compatibility.

As "undefined" is an implicit type, the developer can not add a description text. For "undefined" compatibility the extension repository client will provide two possibilities:
\begin{enumerate}
\item soft-update - upgrading without removing any user data (hence by updating the schema still some data might get lost).
\item hard-update - remove the old version including all user data and reinstall the other version then (possibly the worst attempt).
\end{enumerate}

A corresponding upgrade/downgrade compatibility might look like the following:
\lstset{language=XML} 
\begin{lstlisting}[caption=XML example: upgrade/downgrade compatibility]
...
<ctx:extension version="1.2.3">
    <ctx:upgrade version="1.1" userdata="lossless"/>
    <ctx:upgrade version="1.0" userdata="partial"/>
    <ctx:downgrade version="1.2" userdata="lossless"/>
    <ctx:downgrade version="1.1" userdata="dropUnused"/>
    <ctx:downgrade version="1.0" userdata="partial"/>
</ctx:extension>
...
\end{lstlisting}

When defining such compatibilities, always the most recent version is defined for upgrading and the most oldest version for downgrading.

In above example we have the following definitions:
\begin{tabular}{llp{0.55\textwidth}}\\
\textbf{Rule} & \textbf{Version} & \textbf{Description}\\
upgrade & 1.1 & All versions 1.1 <= X < 1.2.3 can be updated without data loss. \\
upgrade & 1.0 & All versions 1.0 <= X < 1.1 can be updated with partial data loss. \\
downgrade & 1.2 & To all versions 1.2 <= X < 1.2.3 can be downgraded without data loss. \\
downgrade & 1.1 & To all versions 1.1 <= X < 1.2 can be downgraded and unused data will be removed. \\
downgrade & 1.0 & To all versions 1.0 <= X < 1.1 can be downgraded with loss of user data. \\
\end{tabular}
The ordering of information in the xml tree is of no relevance as all items have to be sorted by the client application and sorted. Then the ranges will be defined and the proper information has to be retrieved.

\subsection[sec:files information file]{FILES - package contents information file}
The package contents information contains a list of all files contained within the package (except for the information file "FILES" itself).
The file is used for validation of the installation and especially needed when creating a signed package file to validate the integrity of the installation.
When creating an unsigned package, the file is not required and therefore may be omitted.

The file syntax is simple.
Every line is made up of a sha1 hash followed by a space character and the full file path specifying the destination.

\lstset{language=XML} 
\begin{lstlisting}[caption=FILES example: package contents information file]
...
4543a0fe56f27277948469c15284ac25 CONTAO/system/modules/myExtension/Class1.php
a692ffb1c6026092d916872b2114941c CONTAO/system/modules/myExtension/Class2.php
2709656afab36a8f635206c6956cbf53 CONTAO/system/modules/myExtension/templates/class1\_template.tpl
3fced4c880f491dce98c116266361f1c CONTAO/system/modules/myExtension/templates/class2\_template.tpl
b52f2d57d10c4f7ee67a7eb9615d5d24 LGPL
2e7781b7b08ab489e08526f43adf5b0c PACKAGE
...
\end{lstlisting}

\section[sec:package relationship name spaces]{Package relationship name spaces}

\subsection[sec:package relationship name spaces php features]{PHP features}

\todo{this section is currently undocumented}

\subsection[sec:package relationship name spaces pear packages]{PEAR features}

\todo{this section is currently undocumented}

\subsection[sec:package relationship name spaces database drivers]{Database driver}

\todo{this section is currently undocumented}


% Erstellung eines Index und Abkürzungsverzeichnisses aktivieren ---------------
\makeindex
\makenomenclature
\makeglossaries

% Kopf- und Fußzeilen, Seitenränder etc. ---------------------------------------
% Zeilenabstand 1,5 Zeilen -----------------------------------------------------
%\onehalfspacing
\spacing{1.0}

% Seitenränder -----------------------------------------------------------------
\setlength{\topskip}{\ht\strutbox} % behebt Warnung von geometry
\geometry{paper=a4paper}

% Kopf- und Fußzeilen ----------------------------------------------------------
\pagestyle{scrheadings}
% Kopf- und Fußzeile auch auf Kapitelanfangsseiten
\renewcommand*{\chapterpagestyle}{scrheadings} 
% Schriftform der Kopfzeile
\renewcommand{\headfont}{\normalfont}

% Kopfzeile
\ihead{\large{\textsc{\titel}}\\ \small{\untertitel} \\[2ex] \textit{\headmark}}
\chead{}
%\ohead{\includegraphics[scale=0.075]{\logo}}
\setlength{\headheight}{21mm} % Höhe der Kopfzeile
% Kopfzeile über den Text hinaus verbreitern
%\setheadwidth[0pt]{textwithmarginpar} 
\setheadsepline[text]{0.4pt} % Trennlinie unter Kopfzeile

% Fußzeile
\ifoot{\copyright\ \autor}
\cfoot{}
\ofoot{\pagemark}

% sonstige typographische Einstellungen ----------------------------------------

% erzeugt ein wenig mehr Platz hinter einem Punkt
\frenchspacing 

% Schusterjungen und Hurenkinder vermeiden
\clubpenalty = 10000
\widowpenalty = 10000 
\displaywidowpenalty = 10000

% Quellcode-Ausgabe formatieren
\lstset{numbers=left, numberstyle=\tiny, numbersep=5pt, breaklines=true}
\lstset{emph={square}, emphstyle=\color{red}, emph={[2]root,base}, emphstyle={[2]\color{blue}}}

% Fußnoten fortlaufend durchnummerieren
\counterwithout{footnote}{chapter}

%% ... und nimm alle 4 Ebenen in das Inhaltsverzeichnis auf.
\setcounter{tocdepth}{4}

% eigene Definitionen für Silbentrennung
% Trennvorschläge im Text werden mit \" angegeben
% untrennbare Wörter und Ausnahmen von der normalen Trennung können in dieser
% Datei mittels \hyphenation definiert werden


% eigene LaTeX-Befehle
% Eigene Befehle und typographische Auszeichnungen für diese

% einfaches Wechseln der Schrift, z.B.: \changefont{cmss}{sbc}{n}
\newcommand{\changefont}[3]{\fontfamily{#1} \fontseries{#2} \fontshape{#3} \selectfont}

% Abkürzungen mit korrektem Leerraum 
\newcommand{\ua}{\mbox{u.\,a.\ }}
\newcommand{\zB}{\mbox{z.\,B.\ }}
\newcommand{\dahe}{\mbox{d.\,h.\ }}
\newcommand{\Vgl}{Vgl.\ }
\newcommand{\bzw}{bzw.\ }
\newcommand{\evtl}{evtl.\ }
\newcommand{\ie}{\mbox{i.\,e.\ }}

\newcommand{\illustration}[1]{illustration~\ref{fig:#1}}

\newcommand{\bs}{$\backslash$}

\newcommand{\sectiontitle}{}
\let\TeXsection\section
\renewcommand{\section}[2][\empty]{\ifthenelse{\equal{#1}{\empty}}{\TeXsection{#2}\renewcommand{\sectiontitle}{#2}}{\TeXsection{#2}\label{#1}\renewcommand{\sectiontitle}{#1}}}

\newcommand{\subsectiontitle}{}
\let\TeXsubsection\subsection
\renewcommand{\subsection}[2][\empty]{\ifthenelse{\equal{#1}{\empty}}{\TeXsubsection{#2}\renewcommand{\subsectiontitle}{#2}}{\TeXsubsection{#2}\label{#1}\renewcommand{\subsectiontitle}{#1}}}

\let\TeXsubsubsection\subsubsection
\renewcommand{\subsubsection}[2][\empty]{\ifthenelse{\equal{#1}{\empty}}{\TeXsubsubsection{#2}}{\TeXsubsubsection{#2}\label{#1}}}

\newcommand{\supref}[1]{\textsuperscript{\ref{#1}}}

% erzeugt ein Listenelement mit fetter Überschrift 
\newcommand{\itemd}[2]{\item{\textbf{#1}}\\{#2}}

% einige Befehle zum Zitieren --------------------------------------------------
\newcommand{\Zitat}[2][\empty]{\ifthenelse{\equal{#1}{\empty}}{\citep{#2}}{\citep[#1]{#2}}}

% zum Ausgeben von Autoren
\newcommand{\AutorName}[1]{\textsc{#1}}
\newcommand{\Autor}[1]{\AutorName{\citeauthor{#1}}}

% verschiedene Befehle um Wörter semantisch auszuzeichnen ----------------------
\newcommand{\newDef}[1]{\textbf{#1}}
\newcommand{\techTerm}[1]{\textit{#1}}

\newcommand{\DisplInput}[1]{\texttt{#1}}
\newcommand{\DisplCode}[1]{\texttt{#1}}
\newcommand{\DisplFile}[1]{\texttt{#1}}

\newcommand{\Datentyp}[1]{\textsf{#1}}
\newcommand{\XMLElement}[1]{\textsf{#1}}
\newcommand{\Webservice}[1]{\textsf{#1}}

% Deklaration von Hinweistexten
\newcommand{\hint}[1]{
\begin{tabular}{w{1.5cm}p{12cm}}\\
\scalebox{0.125}{\includegraphics{images/base/Hint_Notification}} & #1
\end{tabular}
}

% Deklaration von Warnungen
\newcommand{\warning}[1]{
\begin{tabular}{w{1.5cm}p{12cm}}\\
\scalebox{0.125}{\includegraphics{images/base/Warning_Notification}} & #1
\end{tabular}
}

% einbinden eines emp Bildes.
\newcommand{\useempcenter}[2]{\begin{figure}[h!]\centering\scalebox{0.8}{\empuse{#1}}\caption{#2}\end{figure}}

\newcommand{\useempfloat}[3]{\begin{wrapfigure}{#1}{0.4\textwidth}\begin{center}\empuse{#2}\end{center}\caption{#3}\end{wrapfigure}}

\newcommand{\useemp}[3][\empty]{\ifthenelse{\equal{#1}{\empty}}{\useempcenter{#2}{#3}}{\useempfloat{#1}{#2}{#3}}}

% Wrapper für useemp für uml Klassendiagramme
\newcommand{\umlclass}[3][\empty]{\ifthenelse{\equal{#1}{\empty}}{\useemp{#2}{#3}}{\useemp[#1]{#2}{#3}}}

% Wrapper um listings einzubinden
\newcommand{\phplisting}[3]{\lstset{language=PHP, basicstyle=\footnotesize, showstringspaces=false, tabsize=2}\lstinputlisting[label=#1,caption=#2]{#3}}
\newcommand{\xmllisting}[3]{\lstset{language=XML, basicstyle=\footnotesize, showstringspaces=false, tabsize=2}\lstinputlisting[label=#1,caption=#2]{#3}}


% Das eigentliche Dokument -----------------------------------------------------
%   Der eigentliche Inhalt des Dokuments beginnt hier. Die einzelnen Seiten
%   und Kapitel werden in eigene Dateien ausgelagert und hier nur inkludiert.
% ------------------------------------------------------------------------------
\begin{document}

% auch subsubsection nummerieren
\setcounter{secnumdepth}{3}
\setcounter{tocdepth}{3}

% Deckblatt und Abstract ohne Seitenzahl
\ofoot{}
\thispagestyle{plain}
\begin{titlepage}

\begin{center}

\huge{\textbf{\titel}}\\[1.5ex]
\LARGE{\textbf{\untertitel}}\\[6ex]
%\includegraphics[scale=0.2]{\deckblattlogo}\\[6ex]

\normalsize
\autor\\[1.2ex]
\copyright\ \jahr\\[9ex]

\end{center}

\singlespacing
\small
\noindent 
\begin{center}
This work is licensed under a \\ \href{http://creativecommons.org/licenses/by-sa/3.0/}{Creative Commons Attribution-ShareAlike 3.0 Unported License}.
\end{center}

A human-readable summary of the Legal Code:

You are free:
\begin{itemize}
\item to Share — to copy, distribute and transmit the work
\item to Remix — to adapt the work
\item to make commercial use of the work
\end{itemize}

Under the following conditions:
\begin{itemize}
\item Attribution — You must attribute the work in the manner specified by the author or licensor (but not in any way that suggests that they endorse you or your use of the work).
\item Share Alike — If you alter, transform, or build upon this work, you may distribute the resulting work only under the same or similar license to this one.
\end{itemize}

Read the license for further details.
\href{http://creativecommons.org/licenses/by-sa/3.0/legalcode}{http://creativecommons.org/licenses/by-sa/3.0/legalcode}

\end{titlepage}

\chapter{Abstract}

This document serves as a discussion and documentation foundation for the new Extension Repository version 3 which shall replace the existing repository for Contao.

This document will be proposed for discussion at Contao Camp 2011.

The first goal of this document is, to identify problems with the current extension repository and provide solutions to them.

To fulfil this goal, the following conditions have to be met:
\begin{itemize}
\item Compatibility to Contao version 2 and version 3 (upcoming).
\item Be more flexible than the current system.
\item (almost) anyone shall be able to run his/her own repository.
\item Installation and uninstallation routines to aid extensions in setup.
\item Standardized package format to ease installation.
\item Improved dependencies.
\item Add possibility to define conflicting extensions.
\end{itemize}

The open discussion on this document is currently taking place in the 
\href{http://www.contao-community.de/showthread.php?22031-Gruppe-neuer-ER-Client-und-ER3}{German Contao Community}

\warning{This document is still a work in progress and should not be used for implementation yet as all parts are still subject to change. Once the document is finalized, it will get a version number and will be released along with some example implementation of the services.}
\ofoot{\pagemark}


% arabische Seitenzahlen im Hauptteil ------------------------------------------
\clearpage
\pagenumbering{arabic}
\tableofcontents % Inhaltsverzeichnis

% UML support
\begin{empfile}
\begin{empcmds}
input metauml;
\end{empcmds}
% die Inhaltskapitel werden in "Content.tex" inkludiert -------------------------
% Hier können die einzelnen Kapitel inkludiert werden. Sie müssen in den 
% entsprechenden .TEX-Dateien vorliegen. Die Dateinamen können natürlich 
% angepasst werden.

\chapter{Changelog}

\newcommand{\ChangeLogEntry}[3]{\subparagraph{Changes by #1 on #2} \begin{itemize} #3 \end{itemize} }

\tiny
\ChangeLogEntry{Tristan Lins}{08.09.2011}{
	\item Initial draft containing document description and package definition.
}
\ChangeLogEntry{Tristan Lins}{09.09.2011}{
	\item added PACKAGE XML information file.
	\item small adjustments in PACKAGE INI Format.
	\item elaborated usage of versions with an extension component.
}
\ChangeLogEntry{Tristan Lins}{10.09.2011}{
	\item Added first description of ER3 Repository, ER3 Repository Manager and ER3 Portal.
}
\ChangeLogEntry{Tristan Lins}{11.09.2011}{
	\item Removed INI Meta file specification.
	\item Reworked PACKAGE information file.
	\item Reworked XML schema.
	\item Reworked XML example.
}
\ChangeLogEntry{Tristan Lins}{12.09.2011}{
	\item Reworked XML Schema and example.
	\item added "versioning of dependencies" and "up/downgrade compatibility" to PACKAGE information.
	\item Removed description of license file and EULA.
}
\ChangeLogEntry{Christian Schiffler}{13.09.2011-14.09.2011}{
	\item Transition to \LaTeX{}, 
	\item translation to English and major rephrasing of all sections
	\item removed \techTerm{lowest\_bound} and \techTerm{highest\_bound} from \DisplFile{extension.xsd}.
}
\ChangeLogEntry{Tristan Lins}{15.09.2011}{
	\item Reworked Repository Portal\supref{sec:repository portal} section
}
\ChangeLogEntry{Christian Schiffler}{16.09.2011}{
	\item finished translation to English (mainly repository\supref{sec:repository} and repository manager \supref{sec:repository manager}).
}
\ChangeLogEntry{Tristan Lins}{19.09.2011}{
	\item Removed package scrhack, it generates the firstpage with "stylistings".
	\item Added first description of "Installation"\supref{sec:installation}, "Installation Process"\supref{sec:installation process} and "Installer"\supref{sec:installer interfaces and classes}.
}
\ChangeLogEntry{Tristan Lins}{25.09.2011}{
	\item Rework complete chapter "Installation"
	\item Added section "ER3 Manager"\supref{sec:er3 manager}.
	\item Added section "Contao Installation Manager"\supref{sec:contao installation manager}.
	\item Rename section "Installer" to "Installer interfaces and classes"\supref{sec:installer interfaces and classes}.
	\item Added subsection "InstallerInformation"\supref{sec:installerinformation}.
}
\ChangeLogEntry{Tristan Lins}{03.10.2011}{
	\item Added subsection "Extension Builder"\supref{sec:extension builder} to section "Repository Portal"\supref{sec:repository portal}
	\item Added section "Name spaces"\supref{sec:namespaces}.
}
\ChangeLogEntry{Tristan Lins}{12.10.2011}{
	\item Added subsection "Comments and rating"\supref{sec:comments and rating} to section "Repository Portal"\supref{sec:repository portal}
}
\ChangeLogEntry{Tristan Lins}{16.10.2011}{
	\item Added encrypted flag to "PACKAGE information file schema"\supref{sec:package information file schema}
	\item Added chapter Licensing
	\item Added appendix "Unencrypted package checksum calculation (upc.php)"\supref{sec:unencrypted package checksum calculation}
	\item Added appendix "Package encryption (encrypt.php)"\supref{sec:package encryption implementation}
	\item Added appendix "Installation id calculation (installationid.php)"\supref{sec:unencrypted package checksum calculation}
	\item Added example "Hello World!" extension package
}
\ChangeLogEntry{Tristan Lins}{25.10.2011}{
	\item Added chapter Certificates\supref{sec:certificates}
	\item Added chapter Extension store\supref{sec:extension store}
	\item Update external hyper links, now they works correctly.
}
\ChangeLogEntry{Tristan Lins}{26.10.2011}{
	\item Added chapter Missing topics\supref{sec:missing topics}
}
\ChangeLogEntry{Tristan Lins}{05.11.2011}{
	\item Added implementation of Version class
	\item Added unit test for Version class
}
\ChangeLogEntry{Tristan Lins}{13.11.2011}{
	\item Added chapter groups and extension management access
	\item Added chapter forks and successors
	\item Added chapter lifetime
}

\ChangeLogEntry{Christian Schiffler}{16.05.2012}{
	\item fixed typos
}
\normalsize
\label{sec:missing topics} 
\chapter{Missing topics}

Topics that needs to be added!

\paragraph{I18N Packages}
A description how internationalisation of packages is done.

\paragraph{UI Framework}
All the web interfaces are modular and able to used without each other. This modularity require a modular ui framework.


\chapter{Packages}

\section[sec:versioning]{Versioning}
The versioning is based on the classic version numbering scheme with extensions.

A version number is combined like the following.

\begin{lstlisting}[caption=Version number in EBNF]
DOT            = <the dot (".") character>
UPALPHA        = <any US-ASCII upper case letter "A".."Z">
LOALPHA        = <any US-ASCII lower case letter "a".."z">
DIGIT          = <any US-ASCII digit "0".."9">
ALPHA          = UPALPHA | LOALPHA
ALNUM          = ALPHA | DIGIT
MAJOR          = DIGIT
MINOR          = DIGIT
REVISION       = DIGIT
EXTENSION      = <"-"> *ALPHA *DIGIT
BUILD          = *DIGIT
HEADBUILD      = <"head"> [DOT *ALNUM]
REGULARBUILD   = <"alpha" DIGIT | "beta" DIGIT | "rc" DIGIT | "stable"> [DOT BUILD]
RELEASE        = HEADBUILD | REGULARBUILD

VERSION = MAJOR DOT MINOR [DOT REVISION] [EXTENSION] DOT RELEASE
\end{lstlisting}

The only mandatory fields in the version are major, minor and the release. 
The release number is combined of either the word "head" in conjunction with an alpha numeric revision suffix which refers to unstable releases or one of the words "alpha", "beta", "rc" or "stable" suffixed with one digit. The latter one can also be accompanied with a build number.
The revision number (also known as "maintenance" or "patch level") is entirely optional, as it is not used in most cases.

The optional extension component is mostly useful for personal repositories, which serve altered versions of certain extensions. This is known from derivative works like the Ubuntu project, which is releasing altered debian packages with a suffix "ubuntu1" or similar.

Please note that developers are encouraged to provide adjustments to upstream maintainers in favor of providing own versions of packages though to keep the forking of extensions to a minimum.

\subsection[sec:versioning real word examples]{A real world versioning example}

Let's imagine an agency has altered an extension for ACME Inc. and is providing this extension via their own repository.
To eliminate problems with the vanilla extension in the main repository, the agency appends the extension "acme" to the version and results in the version
\DisplFile{1.2.3-acme1}.

Using this scheme one can prevent unwanted upgrading (see \ref{sec:versioncompatibility}).

In contrast to the other components of the version, the build number is the only one that can not be influenced by the developer. The build number is incremented by one for every package being uploaded into the repository but not when the description or the properties are changing. This is in contrast to the old ER, which incremented the build number upon every change to the extension.
\todo{How shall this work? version bumping can only be done upon "building" of a package but the upload is much later in the work flow.}


\section[sec:versioncompatibility]{Version compatibility}
Thanks to the \techTerm{extension} component of the version, it is possible to tag modified variants of an extension. These variants are in most cases incompatible to newer (vanilla) releases of the same extension.

To address these potentional problems, the following version compatibility rules are in place.

Versions are compatible if, and only if, any of the following conditions is met:
\begin{itemize}
  \item ...the alphabetic representation of the \techTerm{extension} component is identical.
  \item ...the extension component is omitted on both versions.
\end{itemize}

\small
\begin{longtable}{|p{0.25\textwidth}|p{0.05\textwidth}|p{0.25\textwidth}|}
\caption{Version compatibility examples (compatible: "="; incompatible: "!")} \\
\hline
\label{tab:versioncompatexamples}
\textbf{Version 1} & \textbf{=/!} & \textbf{Version 2} \\
\hline
 1.2.3 & = & 4.5.6
\\ \hline
 1.2-extA.3 & = & 4.5-extA.6
\\ \hline
 1.2-extB.3 & ! & 4.5-extC.6
\\ \hline
 1.2-extD3.4 & = & 5.6-extD7.8
\\ \hline
 1.2-extE3.4 & ! & 5.6-extF7.8
\\ \hline
\end{longtable}
\normalsize

Please note that this only refers to the compatibility in respect to the version and that the compatibility can be further restricted by the extension meta file as described in \ref{sec:package information file}.

\section[sec:version comparison]{Version comparison}

When comparing compatible versions (according to the rules defined in \ref{sec:versioncompatibility}), one has to compare versions with extension component with a different approach than those without extension component.

\subsection[sec:version comparison without extension]{Version comparison without extension component}
A version without extension component will be matched according to the usual mechanism.

A given version A is bigger (aka. newer) than version B if, and only if, one of the following conditions is true:
\begin{itemize}
\item A.major>B.major
\item A.major=B.major \& A.minor>B.minor
\item A.major=B.major \& A.minor=B.minor \& A.revision>B.revision
\item A.major=B.major \& A.minor=B.minor \& A.revision=B.revision \& A.build>B.build
\end{itemize}

A given version A ist identical to version B if, and only if:
\begin{itemize}
\item A.major=B.major \& A.minor=B.minor \& A.revision=B.revision \& A.build=B.build
\end{itemize}

\subsection[sec:version comparison with extension]{Version comparison with extension component}

If an \techTerm{extension} component is present, the following additional rules are to be applied in addition to the rules mentioned in \ref{sec:version comparison without extension}.

A given version A is bigger (aka. newer) than version B if, and only if:
\begin{itemize}
\item A.extensionA = B.extensionA \& A.extensionN > B.extensionN
\end{itemize}

A given version A ist identical to version B if, and only if:
\begin{itemize}
\item A.extensionA = B.extensionA \& A.extensionN = B.extensionN
\end{itemize}

At this, one must pay attention to the fact that the alphabetical representation and the numeric representation has to be evaluated separately.

If the alphabetical representation (extensionA) is different for two versions, they have to be regarded as incompatible according to \ref{sec:versioncompatibility}. When no numeric representation (extensionN) is present in a version with extension, the version has to be evaluated as if a numeric part has been set with the value of "1".

\small
\begin{longtable}{|p{0.25\textwidth}|p{0.25\textwidth}|p{0.25\textwidth}|}
\caption{Version comparison example} \\
\hline
\label{tab:versioncompareexamples}
\textbf{field} & \textbf{Version 1} & \textbf{Version 2} \\
\hline
full version & 1.2-acme3.4 & 1.2-acme4.1
\\ \hline
major & 1 & 1
\\ \hline
minor & 2 & 2
\\ \hline
extension & acme3 & acme4
\\ \hline
extensionA (alpha) & acme & acme
\\ \hline
extensionN (numeric) & 3 & 4
\\ \hline
\end{longtable}
\normalsize

\section[sec:packet structure]{Package structure}
Extensions are shipped in a package file. Such a package contains meta information, optional package management routines for installation, un-installation and handling upgrades in addition to the application data.

The package file is a zip file with a changed extension to mark it as a package file. This approach is already known \ie from java (jar files).

For the file extension currently the following are considered (on the right hand are the file formats listed that are known to use this file extension):
\\
\begin{tabular}{lp{0.9\textwidth}}\\
.cext & none known so far\\
.cex  & INMOS Transputer Development System Occam User Program, ThumbsPlus File, The Currency Exchanger Rate File, CLAN Output File (Child Language Data Exchange System)\\
.ctp  & American Greetings CreataCard (Broderbund), BestAddress HTML Editor Combined Template File (Multimedia Australia Pty. Ltd.)\\
.ctx  & CTX Chemical File (Gasteiger Group), Chinese Character Input File, Compressed Text, GE Industrial Systems CIMPLICITY Text Version HMI Screen, Alphacam Compiled Text (Planit), Online Course Text (Microsoft Corporation), Pretty Good Privacy (PGP) Ciphertext File (PGP Corporation), Windows Terminal Server INI Backup File (Microsoft Corporation), Visual Basic User Control Binary File, CTRAN/W DEFINE Compressed Data File (Geo-Slope International), Nokia PC Suite Backup Contact File (Nokia)\\
\end{tabular}

A package contains the following files:
\begin{itemize}
\item PACKAGE - package information (required, see \ref{sec:package information file})
\item FILES - list of files contained in the package (required for certification, see \ref{sec:files information file})
\item LICENSE files - the license files for the extension (required according to \ref{sec:package information file})
\item ... any files the developer might need below the virtual root "CONTAO" which refers to the contao installation root.
\item ... any files the developer might need below the virtual root "DATA" for the user data root.
\end{itemize}

\subsection[sec:package information file]{PACKAGE information file}

The XML format is described in the schema file extension.xsd\supref{lst:extension.xsd}.

As usage example serves extension.xml\supref{lst:extension.xml}.

\subsubsection[sec:dependencies conflicts suggestions]{Package relationships - dependencies, conflicts and preferations}

A relationship of an extension to another extension can be defined in form of a \techTerm{dependency}, a \techTerm{conflict} or a \techTerm{preferation}.
Each of these relationships can be towards any version of the mentioned extension or even be narrowed down to certain version numbers.

\subsubsection[sec:versioning of relationships]{Versioning of relationships}

As dependencies and conflicts are often regarding certain versions or a certain range of versions.

In the definition of an relationship, we may use the following attributes:
\begin{itemize}
\item \techTerm{match} which means an exact match on the version (X=Y).
\item \techTerm{lowest} which is implicit inclusive and tells the lower end of the range (X>=Y).
\item \techTerm{highest} which is implicit exclusive and tells the upper end of the range (X<Y).
\end{itemize}
The attributes \techTerm{lowest} and \techTerm{highest} do not have to be used in combination, therefore it is possible to define an open range.
When used in combination, they define a closed range.

The attributes \techTerm{lowest} and \techTerm{highest} must not be used in conjunction with the attribute \techTerm{match}. If used in conjunction with \techTerm{match}, the exact match takes precedence over the other attributes and therefore is always exact.

In the package information file, this might look like the following when defining a conflict for a certain db driver (in this case mysql).
\lstset{language=XML} 
\begin{lstlisting}[caption=XML example: conflict example mysql]
...
<ctx:conflict name="db:mysql">
    <ctx:version match="4"/>
    <ctx:version lowest="5" highest="5.1"/>
    <ctx:version lowest="5.2"/>
</ctx:conflict>
...
\end{lstlisting}

In above example definition we have a conflict for:
\begin{itemize}
\item all version 4 drivers.
\item all versions 5 to 5.1 ( 5<=X<5.1)
\item all versions from 5.2 on upwards (5.2<=X)
\end{itemize}
With above condition we define a conflict for all versions of MySQL except for version 5.1 (we could have done it more simple though).

\subsubsection[sec:handling of versions with extension component]{Handling of versions with extension component}

In general the extension component of a version is ignored as it makes common sense that an altered version of an extension is compatible to it's unaltered counterpart from an external point of view. This behaviour can be altered by defining the extension version in the relationship explicitly. If specified, the match for the version must be performed including the extension part.

\subsubsection[sec:special relationships]{Special relationships}

Package relationships can not only refer other packages but some other name-spaces as well.
This is useful to restrict dependency on some php features, a certain database driver, some PEAR package or something else that has been defined in the specification.

Currently the following name spaces exist.\\
\begin{tabular}{ll}\\
contao         & The Contao core.\\
php:[feature]  & Relates to a certain php feature. See php features in \ref{sec:package relationship name spaces php features}\\
pear:[package] & Relates to a certain pear package. See pear packages in \ref{sec:package relationship name spaces pear packages}.\\
db:[type]      & Relates to a certain database, \ie db:mysql. See database drivers in \ref{sec:package relationship name spaces database drivers}.\\
\end{tabular}

\subsubsection[sec:alternative dependencies]{Alternative dependencies}
Often it is required to have a certain feature available but multiple ways exist to get this feature available when depending on some PHP feature but to also provide a wrapper class as workaround if it is not available.

Instead of defining features within the repository client and therefore having to update the client all the time when a new feature has been defined, we provide the developer with the possibility to define alternative paths by specifying groups of dependencies.

When defining groups of dependencies, the first positive match found is accepted and the group evaluation will be stopped. Note that the order in the XML has no meaning and therefore the attribute \techTerm{preferation} must be used to define which dependency is the most preferable.

\lstset{language=XML} 
\begin{lstlisting}[caption=XML example: alternative relationships]
...
<ctx:group preferation="php:soap nusoap">
<ctx:dependency name="php:soap">
    <ctx:description>I love soap, thats because i need this php extension.</ctx:description>
</ctx:dependency>
<ctx:dependency name="nusoap">
    <ctx:description>If you dont have soap, I pleased with an alternative implementation.</ctx:description>
</ctx:dependency>
</ctx:group>
...
\end{lstlisting}
This example describes, that it depends on the PHP SOAP extension but alternatively can also use the extension "nusoap" if the first is not available.

\subsubsection[sec:upgrade and downgrade compatibility]{Upgrade and downgrade compatibility}
In many cases upgrades and downgrades may cause problems. This is even more the case, when the data schema is changing.
That's the reason why the developer needs a way to define what potential threat an upgrade or downgrade of a certain package implies.

This potentional threat can be categorized in the following sections:

\begin{tabular}{llp{0.55\textwidth}}\\
\textbf{Category} & \textbf{threat factor} & \textbf{Description}\\
lossless  & none & Up- and downgrade can be performed without data loss.\\
dropUnused  & almost none & Up- and downgrade will remove data that will not be in use any more.\\
partial  & might be dangerous & Up- and downgrade might remove user data. Therefore manual intervention is suggested.\\
undefined & dragons live here! &This setting is implicit defined when the developer has not defined a setting at all. The result might be user data loss, in worst case scenario resulting in total data loss.\\
\end{tabular}

Whenever a developer has not defined a compatibility, the following automatic assumptions are made:\\
When upgrading, the compatibility is defaulting to "partial".
When downgrading, the compatibility is defaulting to "undefined".
The extension repository manager will then display a message where the up/downgrade compatibility is shown in addition to the description text of the compatibility.

As "undefined" is an implicit type, the developer can not add a description text. For "undefined" compatibility the extension repository client will provide two possibilities:
\begin{enumerate}
\item soft-update - upgrading without removing any user data (hence by updating the schema still some data might get lost).
\item hard-update - remove the old version including all user data and reinstall the other version then (possibly the worst attempt).
\end{enumerate}

A corresponding upgrade/downgrade compatibility might look like the following:
\lstset{language=XML} 
\begin{lstlisting}[caption=XML example: upgrade/downgrade compatibility]
...
<ctx:extension version="1.2.3">
    <ctx:upgrade version="1.1" userdata="lossless"/>
    <ctx:upgrade version="1.0" userdata="partial"/>
    <ctx:downgrade version="1.2" userdata="lossless"/>
    <ctx:downgrade version="1.1" userdata="dropUnused"/>
    <ctx:downgrade version="1.0" userdata="partial"/>
</ctx:extension>
...
\end{lstlisting}

When defining such compatibilities, always the most recent version is defined for upgrading and the most oldest version for downgrading.

In above example we have the following definitions:
\begin{tabular}{llp{0.55\textwidth}}\\
\textbf{Rule} & \textbf{Version} & \textbf{Description}\\
upgrade & 1.1 & All versions 1.1 <= X < 1.2.3 can be updated without data loss. \\
upgrade & 1.0 & All versions 1.0 <= X < 1.1 can be updated with partial data loss. \\
downgrade & 1.2 & To all versions 1.2 <= X < 1.2.3 can be downgraded without data loss. \\
downgrade & 1.1 & To all versions 1.1 <= X < 1.2 can be downgraded and unused data will be removed. \\
downgrade & 1.0 & To all versions 1.0 <= X < 1.1 can be downgraded with loss of user data. \\
\end{tabular}
The ordering of information in the xml tree is of no relevance as all items have to be sorted by the client application and sorted. Then the ranges will be defined and the proper information has to be retrieved.

\subsection[sec:files information file]{FILES - package contents information file}
The package contents information contains a list of all files contained within the package (except for the information file "FILES" itself).
The file is used for validation of the installation and especially needed when creating a signed package file to validate the integrity of the installation.
When creating an unsigned package, the file is not required and therefore may be omitted.

The file syntax is simple.
Every line is made up of a sha1 hash followed by a space character and the full file path specifying the destination.

\lstset{language=XML} 
\begin{lstlisting}[caption=FILES example: package contents information file]
...
4543a0fe56f27277948469c15284ac25 CONTAO/system/modules/myExtension/Class1.php
a692ffb1c6026092d916872b2114941c CONTAO/system/modules/myExtension/Class2.php
2709656afab36a8f635206c6956cbf53 CONTAO/system/modules/myExtension/templates/class1\_template.tpl
3fced4c880f491dce98c116266361f1c CONTAO/system/modules/myExtension/templates/class2\_template.tpl
b52f2d57d10c4f7ee67a7eb9615d5d24 LGPL
2e7781b7b08ab489e08526f43adf5b0c PACKAGE
...
\end{lstlisting}

\section[sec:package relationship name spaces]{Package relationship name spaces}

\subsection[sec:package relationship name spaces php features]{PHP features}

\todo{this section is currently undocumented}

\subsection[sec:package relationship name spaces pear packages]{PEAR features}

\todo{this section is currently undocumented}

\subsection[sec:package relationship name spaces database drivers]{Database driver}

\todo{this section is currently undocumented}

\section[sec:namespaces]{Namespaces}

Contao 3 extensions should use the PHP 5.3 name space feature. Contao 2 supports PHP prior 5.3 so this section only affects Contao 3 extensions.

\subsection[sec:namespace lock]{Namespace Lock}
\todo{this section needs to be verified against the current beta of Contao 3.}
Extensions can define a name space lock, that means a name space, that is "soft" locked for this extension.
This is a build in feature of the Extension Creator\supref{sec:extension builder} and should prevent multiple use of the same name space.
Soft lock means, that it is not possible to exclusive lock a complete name space, it is only possible to lock one name space level.
For example if a developer lock the name space \DisplFile{myextension\bs} the name space \DisplFile{myextension\bs subnamespace\bs} is not locked.
The developer has to lock every name space he wants to use.

A developer who try to use a locked name space get a warning that its locked from another extension.
Because this only affects the Extension Builder\supref{sec:extension builder}, it is not possible to absolutely prevent multiple use of the same name space.

\section[sec:groups]{Groups and extension management access}

Groups are collections of individuals, for example all members of a company or a group of developers, that are allowed to manage an extension in the repository.

An administrator of an extension can define one or more additional users or groups that are allowed to do perform administrative tasks on an extension.
These tasks include (but are not limited to):
\begin{itemize}
\item add I18N packages (translation maintainer).
\item change the description texts.
\item add/upload new versions of the extension.
\end{itemize}

The nature of the tasks that can be assigned to the groups are subject to the repository server configuration as some task might be added via plugins.
\section[sec:forks]{Forks and successors}

An extension can define, that it is a fork of another existing extension.
The ER3 Portal will link a fork to the original extension and vice versa.

An extension can also define, that it is an successor of another extension.
Extension renaming can be fully traced by the ER3.
\todo{add information where this is to be defined and how. i.e. link to proper section within xml file.}
\section[sec:lifetime]{Lifetime}

An extension is marked as outdated if the extension does not support the current minor release of Contao and the difference of the last supported minor release and the current minor release is more than 2. If the major differs, the difference must not be bigger than 1.

An extension is marked as deprecated if the extension does not support the current minor release of Contao and the difference of the last supported minor release and the current minor release is more than 5. If the major differs, the difference has to be bigger than 1.

Extensions can be deleted by the owner. When an extension gets deleted, it is still stored for 6 month, but when someone tries to install this extension, a warning is shown that this extension will not be available any more.


\chapter{Repository Stack}

In this chapter you find the whole extension repository stack defined.

\begin{landscape}
% I know, this is not really a state diagram but rather entities but hey.
\begin{figure}
\begin{emp}[Repository Stack](20, 20)

StateInfo.iRepositoryDistributionNetwork;
iRepositoryDistributionNetwork.foreColor := (.47, .94, .27);

NoteInfo.iUMLNote;
iUMLNote.foreColor := (0, .9, .9);


State.Portal("Repository Portal")();
ENote.NotePortal(iUMLNote)(
	"Public browsing interface",
	"- Public registration",
	"- Public profile management",
	"- Social networking");

State.RepManager("Repository Manager")();
ENote.NoteRepManager(iUMLNote)(
	"Authorization required management interface",
	"- User/Account management",
	"- Package builder",
	"- Package management",
	"- Administration");

State.RepMaster("Repository (Master)")();
ENote.NoteRepMaster(iUMLNote)("Package distribution network");

State.RepSlaveA("Repository (Slave1)")();
State.RepSlaveB("Repository (Slave2)")();
State.RepSlaveC("Repository (Slave3)")();

% this is ugly as hell but I was unable to position the notes more properly using
% same.left(NotePortal, NoteRepManager, NoteRepMaster);
% it always threw an error, FIX this.
leftToRight(95)(Portal, NotePortal);
leftToRight(90)(RepManager, NoteRepManager);
leftToRight(150)(RepMaster, NoteRepMaster);

% combine distribution network
leftToRight(20)(RepSlaveA, RepSlaveB);
topToBottom(20)(RepMaster, RepSlaveA);
EState.RepositoryDistributionNetwork(iRepositoryDistributionNetwork)("Package distribution network")(RepMaster, RepSlaveA, RepSlaveB);

% positioning of globals
topToBottom(20)(Portal, RepManager, RepositoryDistributionNetwork);

State.RepositoryStack("Repository Stack")(Portal, RepManager, RepositoryDistributionNetwork);

drawObjects(RepositoryStack, NotePortal, NoteRepManager, NoteRepMaster);

clink(transition)(RepMaster, RepSlaveA);
clink(transition)(RepMaster, RepSlaveB);

clink(transition)(Portal, RepManager);
clink(transition)(RepManager, RepMaster);


clink(dashedLink)(Portal, NotePortal);
clink(dashedLink)(RepManager, NoteRepManager);
clink(dashedLink)(RepMaster, NoteRepMaster);

\end{emp}
\caption{Repository Stack}
\end{figure}

\end{landscape}

\section[sec:repository]{Repository}
Definition:
The repository is a package storage system that serves package information and packages.
To manage packages within the system, a REST API has been defined. Therefore the repository does not provide a frontend on it's own (aside from the setup interface).

The repository can therefore be compared to a web server providing editing possibilities in addition to information retrival.
For the sake of simplicity, it is not basing upon oversized solutions like WebDAV.

\subsection[sec:repository package storage system]{Package storage system layout}
The package storage system contains all packages in a hierarchical directory tree.

\dirtree{%
.1 ROOT.
.2 <firstchar>.
.3 <extensionname>.
.4 <extensionname>-<version>.<file-extension>.
}
Where the pathes are within the installation root.

The result might be something like the following:
\dirtree{%
.1 ROOT.
.2 a.
.3 anextension.
.4 anextension-1.0.0.cext.
.2 m.
.3 myextension.
.4 myextension-0.1.0.cext.
}

\hint{The nature of this directory tree is only of interest to the implementor of a repository server, as the client will never have to know about this layout, as when downloading a package, the url will be retrieved via the REST API and not calculated on client side.}

\subsection[sec:repository setup interface]{Setup interface}

The setup interface is the only web interface provided by the repository. It is used for installation, upgrading, maintenance and setup of access control.

\subsection[sec:repository rest api]{REST API}

\hint{please note that this section is currently in heavy development and that most details of the REST API have yet to be defined. We have not yet decided how the output data shall be structured and even the result values of endpoints like push, put and delete are yet to be defined.}

The REST API serves the purpose of information retrieval (may be used without authentication) and to manage the repository (requires authentification).

The REST API is implemented as HTTP service and can be reached via requests to the "\DisplFile{/rest}" URL of a web server.

Calls to the API must be either GET or POST requests. For details which method has to be used, please examine the documentation of the commands.

\subsubsection[sec:repository rest api url]{REST API URL}

A REST request URL is made up of the command, optionally followed by a variable amount of parameters. The final element is an optional switching of the output format.

The general syntax of a REST request is made up as defined in the following listing.
\begin{lstlisting}[caption=REST API URL in EBNF]
SLASH          = <the dot ("/") character>
UPALPHA        = <any US-ASCII uppercase letter "A".."Z">
LOALPHA        = <any US-ASCII lowercase letter "a".."z">
DIGIT          = <any US-ASCII digit "0".."9">
ALPHA          = UPALPHA | LOALPHA
ALNUM          = ALPHA | DIGIT
FORMAT         = <"json"> | <"xml">
ENDPOINT       = *ALNUM
COMMAND        = *ALNUM
PARAMNAME      = *ALNUM
PARAMVALUE     = *ALNUM
PARAMETER      = (PARAMNAME SLASH PARAMVALUE)
PARAMETERS     = *(SLASH PARAMETER) | PARAMVALUE
URL            = <"rest"> ENDPOINT SLASH COMMAND [PARAMETERS] [SLASH FORMAT]
\end{lstlisting}

Please pay attention to the fact, that parameters always make up a name-value pair, except for commands where only one parameter is allowed (\ie in the delete\supref{sec:repository rest api endpoint delete} endpoint). When this is the case, this very single parameter consists only of the parameter value and therefore is nameless.

\subsubsection[sec:repository rest api output format]{requesting a custom output format}

The requested output format may be specified by appending an suffix to the url.

Currently supported output formats are \DisplCode{json} and \DisplCode{xml}. For details of the resulting output, please refer to the documentation in Section \todo{link to output format sections here.}.

\subsubsection[sec:repository rest api endpoint list]{REST endpoint "list"}

The list endpoints provides a way to request information of extensions from the repository.

\paragraph{command "recent"}

Via the \DisplCode{recent} command one can acquire a list of the packages that have been recently updated.

Parameters:\\
\begin{tabular}{lp{0.9\textwidth}}\\
filter & Filtering of package names. If omitted, this parameter is implicitly empty and therefore no filtering is applied.\\ 
offset & The offset for paginated output. If omitted this parameter is implicitly set to 0. \\ 
limit & The amount of entries to be returned per page. If omitted, this parameter is implicitly set to 20. The maximum is 100 to ease server load.
\end{tabular}

\vspace{.5cm}
\todo{find some better way to present examples}
Example 1: list recent extensions and filter by category "plugin".\\
\DisplCode{/rest/list/recent/filter/plugin}\vspace{.5cm}

Example 2: list recent extensions and skip first 10 entries.\\
\DisplCode{/rest/list/recent/offset/10}\vspace{.5cm}

Example 3: list 10 most recent extensions.\\
\DisplCode{/rest/list/recent/limit/10}\vspace{.5cm}

\paragraph{command "categories"}

Via the \DisplCode{categories} command one can acquire the list of the package categories that are contained in the repository.

Parameters:\\
\begin{tabular}{lp{0.9\textwidth}}\\
offset & The offset for paginated output. If omitted this parameter is implicitly set to 0. \\ 
limit & The amount of entries to be returned per page. If omitted, this parameter is implicitly set to 20. The maximum is 100 to ease server load.
\end{tabular}

\paragraph{command "tags"}

Via the \DisplCode{tags} command one can acquire the list of all tags that are attached to packages in the repository.

Parameters:\\
\begin{tabular}{lp{0.9\textwidth}}\\
offset & The offset for paginated output. If omitted this parameter is implicitly set to 0. \\ 
limit & The amount of entries to be returned per page. If omitted, this parameter is implicitly set to 20. The maximum is 100 to ease server load.
\end{tabular}

\paragraph{command "search"}

Via the \DisplCode{search} command one can acquire a list of packages that match the given search criteria.

Parameters:\\
\begin{tabular}{lp{0.9\textwidth}}\\
keyword & Match a certain keyword. The keyword may be contained in the description or in some other text information anywhere in the extension.\\
filter & Filtering of package names. If omitted, this parameter is implicitly empty and therefore no filtering is applied.\\ 
offset & The offset for paginated output. If omitted this parameter is implicitly set to 0. \\ 
limit & The amount of entries to be returned per page. If omitted, this parameter is implicitly set to 20. The maximum is 100 to ease server load.
\end{tabular}

\subsubsection[sec:repository rest api endpoint put]{REST endpoint "put"}

Via the \DisplCode{put} endpoint, a package is uploaded into the repository. The package is transferred in the HTTP POST data.

\subsubsection[sec:repository rest api endpoint push]{REST endpoint "push"}

Via the \DisplCode{push} endpoint, a package is uploaded into the repository. In contrast to the \DisplCode{put}\supref{sec:repository rest api endpoint put} endpoint, only an url from where the package can be retrieved is transferred in the HTTP POST data. The repository will then fetch the package from that url.

\subsubsection[sec:repository rest api endpoint delete]{REST endpoint "delete"}

Using the delete endpoint, one can remove either one package or all packages of an extension from the archive.

The availability of this endpoint will be most likely restricted on public repositories (like the one on contao.org) as the necessity to delete releases or even complete extensions should be very rare and therefore will be restricted to repository administrators to prevent abuse.

\paragraph{command "release"}
Remove a single release from the repository. This command only takes one parameter, specifying the exact release name in the form \DisplFile{<name>-<version>} (this is in fact the full package name without extension). The file will then be removed from the database.

\warning{You should pay special attention that when removing a release from the repository this will not remove the version from Contao installations where it already has been installed and therefore might leave those installations in a state where the extension can not be properly upgraded or downgraded as the relevant information has been lost. In addition when removing a release from the repository, it may take some time until all slave servers catch up and are back in sync.}

\paragraph{command "extension"}
Remove an extension including all of it's releases from the repository. This command only takes one parameter, specifying the extension name in the form \DisplFile{<name>} (this is the name of a package file without version and extension). The extension will then be removed from the database along with all of it's releases.

\warning{You should pay special attention that when removing an extension from the repository this will not remove the version from Contao installations where it already has been installed. In addition when removing an extension from the repository, it may take some time until all slave servers catch up and are back in sync.}

\subsubsection[sec:repository rest api endpoint settings]{REST endpoint "settings"}

\todo{the original information from the google doc is very vague, we mus elaborate this to some usable state.}
/rest/settings/name/version

returns the settings for an extension in the repository (category, published state, access restrictions)

\subsubsection[sec:repository rest api endpoint configure]{REST endpoint "configure"}

\todo{the original information from the google doc is very vague, we mus elaborate this to some usable state.}
/rest/configure/name/version

Update the settings of an extension. The new data will be provided via POST data.

\subsubsection[sec:repository rest api authentication]{Authentication}

To authenticate against the REST API, the user can choose between two methods.
\begin{enumerate}
\item authentication via username and password.
\item authentication via OAuth and an auth key.
\end{enumerate}
\todo{translation}

All requests can also be restricted to certain IPs via server configuration. Refer to the server configuration manual.\todo{this manual of course still has to be written. :)}

\section[sec:repository manager]{Repository Manager}

The repository manager is a web interface and the most intuitive form of managing a repository. The manager provides all basic functions of an repository and is part of the Portal\supref{sec:repository portal} but can also be used separately.

By combining the repository\supref{sec:repository package storage system} with the repository manager, one can build a simple to manage repository system. In fact this combination is in essence like the repository in version 2 on contao.org (but will supersede it in features and usability) from developer point of view. Extension packages will get built within the manager and published onto the repository via the REST API.

An installation that consists of this minimal installation is mainly of interest for agencies that want to run a private repository and do not intend to provide a public portal.

\section[sec:repository portal]{Repository Portal}

The portal is a public interface. It links external services to provide additional informations for extensions, that are available on the world wide web.

First it is a repository browser. Users can browse and search extensions and show details of them.
Secondary it is a information aggregator to display related informations.

A package contains links to other sources, see \ref{sec:package information file}.
\lstset{language=XML} 
\begin{lstlisting}[caption=XML example: links to other sources]
...
<ctx:link rel="homepage" href="http://www.contao.org">
    <ctx:description>Official project homepage.</ctx:description>
</ctx:link>
<ctx:link rel="manual" href="http://www.contao.org/documentation.html" language="en">
    <ctx:description>Official User-Manual.</ctx:description>
</ctx:link>
<ctx:link rel="help" href="http://en.contaowiki.org" language="en">
    <ctx:description>Inofficial Wiki.</ctx:description>
</ctx:link>
<ctx:link rel="video" href="http://www.youtube.com/watch?feature=player_embedded&amp;v=8HkQVmfjjYc">
    <ctx:description>Introduction into TypoLIGHT/Contao.</ctx:description>
</ctx:link>
<ctx:link rel="forum" href="http://www.contao-community.de" language="de">
    <ctx:description>German contao community.</ctx:description>
</ctx:link>
<ctx:link rel="tracker" href="http://dev.contao.org">
    <ctx:description>Official bug and issue tracker.</ctx:description>
</ctx:link>
<ctx:link rel="scm" href="http://svn.contao.org">
    <ctx:description>Official svn repository.</ctx:description>
</ctx:link>
...
\end{lstlisting}
The portal use this links and try to fetch related informations from these systems. I.e. this can be bugs and issues from a tracking system like Redmine or Github. Or screencast from YouTube. Or tutorials from a MediaWiki.

The Portal is not only an information aggregator. Not only the developer can link external informations to a extensions.
Other users can link ther own informations, like tutorials, screencasts or similar to an extension.
Theses links are stored in the portal, but not in the extension information. This is the social networking character of the Portal.

\subsection[sec:information kinds and system connectors]{Information kinds and system connectors}

The Portal use multiple connectors to get informations from other systems.
First we have to define the kind of informations, we are able to aggregate.

\begin{itemize}
\item Bugs and issues\\
The portal can display bugs and issues to show the current development status and upstream development.
\item Manuals and Tutorials - Wiki informations\\
One of the best ways to manage manuals and tutorials is to use a wiki. The portal can display related wiki articles.
\item Videos and Screencasts\\
Videos and screencasts are another kind of tutorials. The portal can display videos hosted on video streaming platforms like YouTube.
\item Source code\\
If an scm repository is configured, the portal can directly link to the source code.
\end{itemize}

These are just a few possible information kinds, the repository may use.

Next we have to define how the informations are aggregated.
A lot of systems like \href{http://www.mediawiki.org/wiki/API:Main_page/de}{MediaWiki}, \href{http://code.google.com/intl/de-DE/apis/youtube/overview.html}{YouTube} and \href{http://developer.github.com/v3/}{github} provide external APIs.
The Portal use connectors to this systems to get informations.
Because not all theses services are uniq, like MediaWiki the systems have to be configures dynamicaly.
All users are allowed to request adding a service, but only administrators are able to add them finaly into the portal.
Initialy there may be these connectors available:

\small
\begin{longtable}{|p{0.25\textwidth}|c|c|c|c}
\caption{Information system connectors and information kinds} \\
\hline
\label{tab:information system connectors and information kinds}
\textbf{System} & \textbf{Tracking} & \textbf{Wiki} & \textbf{Video} & \textbf{SCM}
\\ \hline
 MediaWiki & & X & &
\\ \hline
 Redmine & X & X & &
\\ \hline
 Github & X & X & & X
\\ \hline
 YouTube & & & X &
\\ \hline
 Vimeo & & & X &
\\ \hline
\end{longtable}
\normalsize

\subsection[sec:comments and rating]{Comments and rating}

Registered users are able to comment and rate extensions.

Every user can only add or edit a single comment to an extension.
This prevents comment flooding.

The rating is a star rating from 1 till 5.
If there are no ratings, the extension get no stars.
With the first rating, the extension gets grey stars in small size (16px).
The stars get more golden and bigger with more ratings.
With the 100th rating, the start are full golden and at the greatest size (32px).

\subsection[sec:extension builder]{Extension Builder}

The Extension Builder is a web interface to simply build extension packages.

\label{sec:extension store} 
\chapter{Extension store}

\todo{Needs work}

The extension store is part of the portal and an "app store" equivalent for contao extensions.


\label{sec:installation} 
\chapter{Installation}

This chapter describe how the Installation goes on and what programming basics are required.

\section[sec:er3 manager]{ER3 Manager}

The ER3 Manager (fka ER Client) is designed for Contao 2. It's a full replacement for the ER2 Client, that is currently in use.
It's a combined manager, that use extensions from the ER2 and the ER3. The ER3 Manager - or Extension Manager - have to manage all installed extension in an existing Contao 2 installation.

\section[sec:contao installation manager]{Contao Installation Manager}

The Contao Installation Manager (short CIM), is an Contao independend tool. It's designed for Contao 3 and should manage a complete Contao installation, including the installed Contao core and the installed extensions.
Summarized it's a combination of the current existing Install Tool and the ER3 Manager.

\paragraph{}
CIM should replace the Contao distribution package.
The idea behind CIM is, that a user only have to load a single PHP file, the CIM Installer.
The CIM Installer load all required CIM files.
Then the CIM provide a wizard to install a complete Contao.
Maybe with a basic template and basic extensions.

\paragraph{}
To make this possible, contao have to managed as a repository package itself.
This means, the CIM can manage the Contao core like any other extension.
In this case, there are no need for special behaviours to manage dependencies to the Contao Core.

\paragraph{}
But CIM is more than an installation tool.
CIM check the environment, if its compatible with Contao and provide solutions to make it compatible
For example activate the FTP file acces (fka SafeModeHack).
CIM manage the Contao system settings.
For example the database connection.
CIM simplified the installation process by automating the installation process.

\paragraph{}
For users, CIM is an absolute simplification of the installation process.
First download the CIM Installer (a small PHP file) and upload to the webspace.
Second call the CIM Installer, the CIM Installer download all required CIM files.
Third make shure, the environment is ready to install.
Fourth load a basic installation (Music Academy).
After CIM finish it work, the installation is ready to use.

\paragraph{}
For web workers/agencies additionaly it automating the installation process.
The automating process can be customized, to mind the custom preferences.


\section[sec:installation process]{Installation Process}

This section describe the flow of the installation process.

\todo{Make some charts to visualize the processes}

\subsection[sec:install process]{Install Process}

\begin{enumerate}
\item Download the package into local package archive
\item Read the installer class from package meta
\item Instanciate the installer
\item Create new installer information object
\item Call the installers \DisplFile{install(..)} method
\item Store the installer information object
\end{enumerate}

\subsection[sec:upgrade process]{Upgrade Process}

\begin{enumerate}
\item Download the new version package into local package archive
\item Read the installer class from package meta
\item Instanciate the installer
\item Load installer information object
\item Call the installers \DisplFile{upgrade(..)} method
\item Store the installer information object
\end{enumerate}

\subsection[sec:downgrade process]{Downgrade Process}

\paragraph{Method 1: Upgrade is possible}

\begin{enumerate}
\item Download the old version package into local package archive
\item Read the installer class from current installed package meta
\item Instanciate the installer of the current installed version
\item Load installer information object
\item Call the installers \DisplFile{downgrade(..)} method
\item Read the installer class from old version package meta
\item Instanciate the installer of the old version
\item Call the installers \DisplFile{upgrade(..)} method
\item Store the installer information object
\end{enumerate}

\paragraph{Method 2: Upgrade not possible - safe method}

\begin{enumerate}
\item Download the old version package into local package archive
\item Read the installer class from current installed package meta
\item Instanciate the installer of the current installed version
\item Load installer information object
\item Call the installers \DisplFile{uninstall(..)} method
\item Read the installer class from old version package meta
\item Instanciate the installer of the old version
\item Call the installers \DisplFile{install(..)} method
\item Store the installer information object
\end{enumerate}

\paragraph{Method 3: Upgrade not possible - hard method}

\begin{enumerate}
\item Download the old version package into local package archive
\item Read the installer class from current installed package meta
\item Instanciate the installer of the current installed version
\item Load installer information object
\item Call the installers \DisplFile{uninstall(..)} method
\item Call the installers \DisplFile{clean(..)} method
\item Read the installer class from old version package meta
\item Instanciate the installer of the old version
\item Call the installers \DisplFile{install(..)} method
\item Store the installer information object
\end{enumerate}

\subsection[sec:reinstall process]{Reinstall Process}

\begin{enumerate}
\item Download the package into local package archive
\item Read the installer class from package meta
\item Instanciate the installer
\item Load installer information object
\item Call the installers \DisplFile{reinstall(..)} method
\item Store the installer information object
\end{enumerate}

\subsection[sec:uninstall process]{Uninstall Process}

\begin{enumerate}
\item Download the package into local package archive
\item Read the installer class from package meta
\item Instanciate the installer
\item Load installer information object
\item Call the installers \DisplFile{uninstall(..)} method
\item Store the installer information object (maybe some DATA files and schema is just remeaning in the system)
\end{enumerate}

\subsection[sec:clean uninstall process]{Clean Uninstall Process}

Same as clean process.

\subsection[sec:clean process]{Clean Process}

\begin{enumerate}
\item Download the package into local package archive
\item Read the installer class from package meta
\item Instanciate the installer
\item Load installer information object
\item If still installed, call the installers \DisplFile{uninstall(..)} method
\item Call the installers \DisplFile{clean(..)} method
\item Remove stored installer information object
\end{enumerate}

\subsection[sec:installation process fallback]{Fallback if package is not available}

This fallback works only for uninstall\supref{sec:uninstall process} process!

If the package is not available, this means the installer can not read wich installer class should be used.
It asked the user to use the basic operation, in this case, the BasicInstaller is used as installer to make the operations.

\section[sec:installer interfaces and classes]{Installer interfaces and classes}

\begin{figure}[h]
\begin{emp}[Installer interfaces and classes](20, 20)
Class.information("InstallationInformation")()(
	"+getVersion() : Version",
	"+getFiles() : array",
	"+getInstalledPackage() : Package",
	"+getTargetPackage() : Package");

Class.installer("Installer")()(
	"+install($objInfo : InstallationInformation, $objVersion : Version) throws InstallerException : void",
	"+upgrade($objInfo : InstallationInformation, $objTargetVersion : Version) throws InstallerException : void",
	"+downgrade($objInfo : InstallationInformation, $objTargetVersion : Version) throws InstallerException : void",
	"+reinstall($objInfo : InstallationInformation) throws InstallerException : void",
	"+uninstall($objInfo : InstallationInformation) throws InstallerException : void",	
	"+clean($objInfo : InstallationInformation) throws InstallerException : void");

Class.basicInstaller("BasicInstaller")()(
	"+removeFiles($objInfo : InstallationInformation) throws InstallerException : void",
	"+removeContaoFiles($objInfo : InstallationInformation) throws InstallerException : void",
	"+removeDataFiles($objInfo : InstallationInformation) throws InstallerException : void",
	"+removeUnusedFiles($objInfo : InstallationInformation) throws InstallerException : void",
	"+removeUnusedContaoFiles($objInfo : InstallationInformation) throws InstallerException : void",
	"+removeUnusedDataFiles($objInfo : InstallationInformation) throws InstallerException : void",
	"+copyFiles($objInfo : InstallationInformation) throws InstallerException : void",
	"+copyContaoFiles($objInfo : InstallationInformation) throws InstallerException : void",
	"+copyUserFiles($objInfo : InstallationInformation) throws InstallerException : void",
	"+removeSchema($objInfo : InstallationInformation) throws InstallerException : void",
	"+safePatchSchema($objInfo : InstallationInformation) throws InstallerException : void",
	"+fullPatchSchema($objInfo : InstallationInformation) throws InstallerException : void");

topToBottom(20)(installer, basicInstaller);
topToBottom(20)(information, installer);

drawObjects(information, installer, basicInstaller);

clink(dashedLink)(information, installer);
clink(transition)(basicInstaller, installer);
\end{emp}
\caption{Installer interfaces and classes}
\end{figure}

\subsection[sec:installerinformation]{InstallationInformation}

The InstallationInformation object, contains all informations about an installation.

\begin{itemize}
\item Installed Version
\item Installed Files
\item Corresponding Schemas
\item The installed package file (if available)
\item The target installation package file (if a new version should be installed or up-/downgraded)
\end{itemize}

The InstallationInformation object is modified from the Installer.
The Installer register all it changes of the extension to the InstallationInformation object, including version changes, file list changes and schema changes.

\subsection[sec:installer]{Installer}

The installer have to provide methods to install, upgrade, downgrade and uninstall an extension.

\DisplCode{install(\$objInfo, \$objVersion)} Installs the extension. It gets the target version to install as argument to make it easer for developers to create multi version installers.

\vspace{.5cm}
\DisplCode{upgrade(\$objInfo, \$objTargetVersion)} Upgrades the extension from one version to another version.

\vspace{.5cm}
\DisplCode{downgrade(\$objInfo, \$objTargetVersion)} Downgrades the extension from one version to another version.

\vspace{.5cm}
\DisplCode{reinstall(\$objInfo)} Reinstall the extension.

\vspace{.5cm}
\DisplCode{uninstall(\$objInfo)} Uninstalls the extension. By definition, this will only uninstall the CONTAO files. All user DATA and schema information and data will be kept.

\vspace{.5cm}
\DisplCode{clean(\$objInfo)} Completely clean the installation, removes all files, including DATA files and schema information.

\vspace{.5cm}
The installation process success if there is no exception thrown.

\subsection[sec:basic installer]{Basic Installer}

The Basic Installer is the default implementation of the installer interface.
If the developer does not define an installer for the extension, this installer is used.
In most cases, if the developer needs to add custom install, upgrade, downgrade or uninstall code, the developer create a supimplementation of this class.

\begin{figure}[h]
\begin{emp}[Basic Installer](20, 20)
\end{emp}
\caption{Basic Installer Implementation}
\end{figure}

% remove methods
\vspace{.5cm}
\DisplCode{removeFiles(\$objInfo)} Removes all CONTAO and DATA files from the installation.

\vspace{.5cm}
\DisplCode{removeContaoFiles(\$objInfo)} Removes all CONTAO files from the installation.

\vspace{.5cm}
\DisplCode{removeDataFiles(\$objInfo)} Removes all DATA files from the installation.

\vspace{.5cm}
\DisplCode{removeUnusedFiles(\$objInfo)} Removes unused CONTAO and DATA files from the installation.

\vspace{.5cm}
\DisplCode{removeUnusedContaoFiles(\$objInfo)} Removes unused CONTAO files from the installation.

\vspace{.5cm}
\DisplCode{removeUnusedDataFiles(\$objInfo)} Removes unused DATA files from the installation.

% copy files methods
\vspace{.5cm}
\DisplCode{copyContaoFiles(\$objInfo)} Copy the files from CONTAO package directory into the installation.

\vspace{.5cm}
\DisplCode{copyDataFiles(\$objInfo)} Copy the files from DATA package directory into the installation.

% schema methods
\vspace{.5cm}
\DisplCode{removeSchema(\$objInfo)} Remove the complete schema and all containing data from the database.

\vspace{.5cm}
\DisplCode{safePatchSchema(\$objInfo)} Add missing or modify existing schema, but does not delete any existing schema information and data.

\vspace{.5cm}
\DisplCode{fullPatchSchema(\$objInfo)} Add missing or modify existing schema, and delete no more existing schema information and data.

\vspace{.5cm}
The basic implementation may look like this. Please note, this is an incomplete, pseudo implementation!

\begin{lstlisting}[caption=Basic Installer Implementation Code]
class BasicInstaller implements Installer
{
	public function install(InstallationInformation \$objInfo, Version \$objVersion)
	{
		$this->copyContaoFiles($objInfo);
		$this->copyDataFiles($objInfo);
		$this->safePatchSchema($objInfo);
	}
	
	public function upgrade(InstallationInformation $objInfo, Version $objToVersion)
	{
		$this->removeUnusedFiles($objInfo);
		$this->safePatchSchema($objInfo);
		$this->copyFiles($objInfo);
	}
	
	public function downgrade(InstallationInformation $objInfo, Version $objToVersion)
	{
		$objVersionCompare = $objInfo->getVersion()->compare($objToVersion);
		
		// check if the versions differ in more than the build number
		if (!$objVersionCompare->isAnotherBuild())
		{
			throw new DowngradeException("Downgrade to another major, minor or revision is not supported.");
		}
		
		$this->removeUnusedFiles($objInfo);
		$this->safePatchSchema($objInfo);
		$this->copyFiles($objInfo);
	}
	
	public function reinstall(InstallationInformation $objInfo, Version $objTargetVersion)
	{
		$this->removeContaoFiles($objInfo);
		$this->safePatchSchema($objInfo);
		$this->copyFiles($objInfo);
	}
	
	public function uninstall(InstallationInformation $objInfo)
	{
		$this->removeContaoFiles($objInfo);
	}
	
	public function clean(InstallationInformation $objInfo)
	{
		$this->removeFiles($objInfo);
		$this->removeSchema($objInfo);
	}
	
	public function removeFiles(InstallationInformation $objInfo)
	{
		$this->removeContaoFiles($objInfo);
		$this->removeDataFiles($objInfo);
	}
	
	public function removeContaoFiles(InstallationInformation $objInfo)
	{
		// remove files from the CONTAO package directory.
		// unregister removed files from $objInfo
	}
	
	public function removeDataFiles(InstallationInformation $objInfo)
	{
		// remove files from the DATA package directory.
		// unregister removed files from $objInfo
	}
	
	public function removeUnusedFiles(InstallationInformation $objInfo)
	{
		$this->removeUnusedContaoFiles($objInfo);
		$this->removeUnusedDataFiles($objInfo);
	}
	
	public function removeUnusedContaoFiles(InstallationInformation $objInfo)
	{
		// remove files from the CONTAO package directory that are unused.
		// unregister removed files from $objInfo
	}
	
	public function removeUnusedDataFiles(InstallationInformation $objInfo)
	{
		// remove files from the DATA package directory that are unused.
		// unregister removed files from $objInfo
	}
	
	public function copyFiles(InstallationInformation $objInfo)
	{
		$this->copyContaoFiles($objInfo);
		$this->copyUserFiles($objInfo);
	}
	
	public function copyContaoFiles(InstallationInformation $objInfo)
	{
		// copy files from the CONTAO package into the installation.
		// register installed files into $objInfo
	}
	
	public function copyUserFiles(InstallationInformation $objInfo)
	{
		// copy files from the DATA package into the installation.
		// register installed files into $objInfo
	}
	
	public function removeSchema(InstallationInformation $objInfo)
	{
		// remove the complete schema
	}
	
	public function safePatchSchema(InstallationInformation $objInfo)
	{
		// add missing schema information and update existing
	}
	
	public function fullPatchSchema(InstallationInformation $objInfo)
	{
		// add missing schema information, update existing and remove no more existing informations.
	}
}
\end{lstlisting}

\section[sec:custom installer code]{Custom Installer}

If a developer want's to add custom code, he can make a subclass of the BasicInstaller.

\begin{lstlisting}[caption=Custom Installer Code Example]
class CustomInstaller extends BasicInstaller
{
	public function upgrade(InstallationInformation $objInfo, Version $objToVersion)
	{
		// add custom pre-upgrade code here
		
		parent::upgrade($objInfo, $objToVersion);
		
		// add custom post-upgrade code here
	}	
}
\end{lstlisting}

Or the developer write his complete own installer. In this case the developer have to implement the Installer interface.


\label{sec:certificates} 
\chapter{Certificates}

The ER3 uses \href{http://en.wikipedia.org/wiki/X.509}{X.509 certificates}.
There are two types of certificates, the developer certificate\supref{sec:developer certificate} for signing packages\supref{sec:signed packages} and the installation certificate\supref{sec:installation certificate} for licensing\supref{sec:licensing}.

\section[sec:developer certificate]{Developer certificate}

Every registered person or company can only have one developer certificate.
The developer certificate is used to sign packages\supref{sec:signed packages}.
The ER3 Manager\supref{sec:er3 manager} can use the signature to validate the integrity and origin.

\section[sec:installation certificate]{Installation certificate}

An installation certificate is only required to use the extension store\supref{sec:extension store}.
It is generated through the ER3 Manager\supref{sec:er3 manager}.
The public key is stored in the Repository Portal\supref{sec:repository portal}, associated to a registered person or company and the installation domain (just for human identification).
A single person or company can register multiple installation certificates.

\paragraph{}
An installation certificate should be unique, this means every installation get its own certificate.
While generating a new installation certificate the ER3 Manager\supref{sec:er3 manager} stored some unique informations with the certificate, the installation domain and the installation path.
If these informations change, the ER3 Manager\supref{sec:er3 manager} ask the user, if a new certificate should be generated.
Normally the answer is no, because in most cases these informations change while put an installation from development to productive, change the provider or something similar.
But if an installation is duplicated to create a new website, for example by using a copy master. This question can prevent installation certificate duplication.
In most cases, this should not be necessary, because an installation certificate is first generated if it is needed. Or in other words, the first time the extension store\supref{sec:extension store} is opened from the ER3 Manager\supref{sec:er3 manager}.

\section[sec:key exchange]{Key exchange}

The public keys are stored in the Repository Portal\supref{sec:repository portal}.
A public key can only be stored by registered persons and companies, this means it is not an open public key storage.
The portal also provide a \href{http://en.wikipedia.org/wiki/Revocation_list}{certification revocation list} for deleted keys.
A key get invalid if the owning person or company account is deleted.

\section[sec:certificate dependencies]{Dependencies}

The php openssl module is required for this functionality.
This module should be available since php 4.

\section[sec:signed packages]{Signed packages}

A developer can sign his packages.
Signed packages provide more security because the user can validate the origin of the package.
All developers should sign there packages, the extension builder\supref{sec:extension builder} normally do this for the developer.
If the ER3 Manager\supref{sec:er3 manager} try to install an unsigned packages it gives a security warning to the user.


\label{sec:licensing} 
\chapter{Licensing}

\todo{This chapter needs rework to use certificates. We are working on it :-)}

There are two types of licenses, the multi license and the single license.
The multi license allow multiple installation of the extension.
The single license allow a single installation of the extension.

A license can be acquired through the "Repository Portal"\supref{sec:repository portal} or from the developer itself.
The "ER3 Manager"\supref{sec:er3 manager} allow to aquire a license directry through the "Repository Portal"\supref{sec:repository portal}.
To use licenses, the "ER3 Manager"\supref{sec:er3 manager} have to be authenticated against the "Repository Portal"\supref{sec:repository portal}.

The license is split into two parts, the license key and the license file.
The license key is a short unique identifier for the license file.
The license key have to be validated against the Repository (not Portal)

\section[sec:license key]{License key}

The license key is a sha256 hash, split into 8 parts with a length of 8 characters, concatenated with a minus character (-).

A license key may look like: \\
\texttt{e3b0c442-98fc1c14-9afbf4c8-996fb924-27ae41e4-649b934c-a495991b-7852b855}

\section[sec:license file]{License file}

The license file contains a lot of informations about the license, stored in INI format:
\begin{itemize}
\item Name of the extension [extension]
\item Lowest extension version [lowest version]
\item Highest extension version [highest version]
\item Name and email of license holder [license holder]
\item Name and email of licensor [licensor]
\item Unencrypted package checksum [checksum]
\item Type of license (single or multi, optional, default is multi) [type]
\end{itemize}
The names in the brackets [ ] are the ini fields.
The fields lowest version and highest version can be combined to a single field version.
This can be an absolut single version or a major or minor version.

\begin{lstlisting}[caption=License file]
[license]
extension=MyContaoExtension
lowest version=1.1
highest version=1
license holder=Tristan Lins <info@infinitysoft.de>
licensor=Christian Schiffler <info@cyberspectrum.de>
checksum=2c26b46b68ffc68ff99b453c1d30413413422d706483bfa0f98a5e886266e7ae
\end{lstlisting}

The license file is never stored in the contao extension. The license file can be optained from the "Repository"\supref{sec:repository} automaticaly by the "ER3 Manager"\supref{sec:er3 manager}. Alternatively the user can download the license file and use it every time, he want to install it (see Standalone license\supref{sec:standalone license}).

\section[sec:unencrypted package checksum]{Unencrypted package checksum}

The unencrypted package checksum (UPC) fulfilled two functions.
The first is to have an unique field in the license file, that cannot be recreated from a user.
The second is to validate the package integrity after decryption.

To calculate the UPC, the folowing steps are required.

Step 1: build a natural ordered list of all files, including the meta files of the package.

Step 2: build a sha256 hash about all concatenaded files, in the order of the list from step 1.

A sample implementation serves upc.php\supref{sec:unencrypted package checksum calculation}.

\section[sec:standalone license]{Standalone license}

A standalone license is used if the "ER3 Manager"\supref{sec:er3 manager} is unable to connect the repository or the repository does not contains the license. There is no different to a normal license, expecting that the license file is not obtained from the "Repository"\supref{sec:repository}. The user have to upload the license file from his computer every time, he want to install or update the extension. This user interaction is necessary because the license file cannot be stored in the contao installation to protect the package from steeling.

For this license usage the "License ipgrade"\supref{sec:license upgrade} is impossible.

\section[sec:encrypted packages]{Encrypted packages}

A licensed package get encrypted to prevent package stealing.
As encryption algorithm the advanced encryption standard (AES) is used.
If a package is encrypted, the \textbf{encrypted} flag have to be set to \textbf{true} and all encrypted files get the extension ".aes".
The encryption only affects files from the "CONTAO" and "DATA" directory.
All meta informations are stored unencrypted.
Encrypted packages requires signing of the packages.
The encryption key can be calculated from the license key and package checksum.
The license checksum (license key without minus) and package checksum are recalculated to bytestreams and then the bytestreams are XOR with each other. The result is the encryption key.
See upc.php\supref{sec:unencrypted package checksum calculation} for implementation example.

The usage of the licence checksum and package checksum as part of the encryption key forces every licensed package to be unique!

A sample implementation serves encrypt.php\supref{sec:package encryption implementation}.

\section[sec:single license]{Single license}

For a single installation there is used a third id to protect the package, the "installation id"\supref{sec:installation id}. For a single license the encryption key is calculated by XOR license checksum, package checksum and installation id. In the license file the "type" field have to set to "single". There are no other changes to the license, only the package use another key for encryption.

\subsection[sec:installation id]{Installation ID}

The installation id is unique for every installation. It is calculated from these factors:
\begin{itemize}
\item Contao Version \todo{Maybe its not the best idea to use Contao Version in the installation id?!}
\item Contao installation path (working directory)
\item Running domain
\end{itemize}

The installation id is a sha256 checksum over the three newline-concatenated values.

A sample implementation serves installationid.php\supref{sec:installation id calculation}.

\section[sec:license upgrade]{Licence upgrade}

As you read the sections above, there is a big problem with these licensing and encryption mechanism.
A license key can only be valid for a single package version.
This is because it is just a hash about the license file and the license file have to contain the checksum of the package.
As a result the "Repository Portal"\supref{sec:repository portal} have to store the information for which version the license is generated.
This is the basic for the license upgrade.

The licensed (encrypted) packages are automatic generated through the "Extension Builder"\supref{sec:extension builder}.
Because this, its possible to automatic generate a new license and package.
If a user try to update a licensed extension and does not own a license for the current version, but for a prior version.
Then the "Repository Portal"\supref{sec:repository portal} automaticale upgrade the license.
In the license file contains all necessary informations for this.
If an upgrade is possible, this means the license includes the current version, a new license file and license key is generated.
The license key is pushed to and stored in the contao installation.
After that, the "ER3 Manager"\supref{sec:er3 manager} can download and upgrade to the new package version.

If the "Extension Builder"\supref{sec:extension builder} is not used by the developer, the developer have to upload an unencrypted package into the repository. The unencrypted package will never be published, but its used to generate the licensed package for every license holder.

% UML support
\end{empfile}

% Seitennummerierung -----------------------------------------------------------
%   Nach dem Hauptteil werden die Seiten in großen römischen Ziffern 
%   nummeriert.
% ------------------------------------------------------------------------------
\pagenumbering{Roman}

% Anhang -----------------------------------------------------------------------
%   Die Inhalte des Anhangs werden analog zu den Kapiteln inkludiert.
%   Dies geschieht in der Datei "Appendix.tex".
% ------------------------------------------------------------------------------
\begin{appendix}
    \clearpage
    \chapter{Appendix}
    \label{sec:Appendix}
    \pagenumbering{roman}
    % Rand der Aufzählungen in Tabellen anpassen
    \setdefaultleftmargin{0em}{}{}{}{}{}
    \section[sec:package information file schema]{PACKAGE information file schema}
\xmllisting{lst:extension.xsd}{PACKAGE information file schema}{../extension.xsd}

\section[sec:package information file example]{PACKAGE information file example}
\xmllisting{lst:extension.xml}{PACKAGE information file example}{../extension.xml}

\section[sec:unencrypted package checksum calculation]{Unencrypted package checksum calculation}
\phplisting{lst:upc.php}{Unencrypted package checksum calculation}{../../src/upc.php}

\section[sec:package encryption implementation]{Package encryption}
\phplisting{lst:encrypt.php}{Package encryption}{../../src/encrypt.php}

\section[sec:installation id calculation]{Installation id calculation}
\phplisting{lst:installationid.php}{Installation id calculation}{../../src/installationid.php}

\end{appendix}

% Abkürzungsverzeichnis --------------------------------------------------------
\nomenclature{API}{Application Programming Interface}
\nomenclature{ARIS}{Architektur integrierter Informationssysteme}
\nomenclature{BPR}{Business Process Reengineering}
\nomenclature{eEPK}{erweiterte Ereignisgesteuerte Prozesskette}
\nomenclature{EPK}{Ereignisgesteuerte Prozesskette}
\nomenclature{JMS}{Java Message Service}
\nomenclature{SDK}{Software Development Kit}
\nomenclature{URI}{Uniform Resource Identifier}
\nomenclature{URL}{Uniform Resource Locator}
\nomenclature{URN}{Uniform Resource Name}
\nomenclature{W3C}{World Wide Web Consortium}
\nomenclature{XML}{Extensible Markup Language}
\nomenclature{XPath}{XML Path Language}
\nomenclature{XSL}{Extensible Stylesheet Language}
\nomenclature{XSLT}{XSL Transformations}

\glossary{name={HOOK}, description={Ein Einsprungpunkt bei dem eine Erweiterung in den Ablauf einer Funktion eingehängt wird.}}
% für korrekte Überschrift in der Kopfzeile
\clearpage\markboth{\nomname}{\nomname} 
\printnomenclature
\label{sec:Glossar}

\listoffigures % Abbildungsverzeichnis
\listoftables % Tabellenverzeichnis
\renewcommand{\lstlistlistingname}{List of Listings}
\lstlistoflistings % Listings-Verzeichnis

% Literaturverzeichnis ---------------------------------------------------------
%   Das Literaturverzeichnis wird aus der BibTeX-Datenbank "Bibliographie.bib"
%   erstellt.
% ------------------------------------------------------------------------------
\bibliography{Bibliographie} % Aufruf: bibtex RFCCC1
\bibliographystyle{natdin} % DIN-Stil des Literaturverzeichnisses

% Index ------------------------------------------------------------------------
%   Zum Erstellen eines Index, die folgende Zeile auskommentieren.
% ------------------------------------------------------------------------------
\printindex
\printglossaries

\end{document}
